\SourceLeaf{260}
\Article{9}
Haec omnia eodem fere modo iam in Disquis.\ Ar.\ pertractata fuerant; quaedam autem illic adiecta fuerant suppressa demonstratione, quam hic explere consultum iudicamus. Annuntiavimus illic, evolutionem valoris quantitatis radicalis ${}^{\beta}\!\surd\theta'$, quae quandoquidem $\theta'$ est quantitas imaginaria, sectionem tum rationis tum anguli in $\beta$ partes requirere videtur, a sola posteriori pendere, prioremque semper ad solam extractionem unius radicis quadratae reduci posse: hoc ita demonstramus.

Designando ut supra quantitatem imaginariam $\surd-1$ per $i$, statuendoque $\theta'=P+iQ$, atque aliquem valorem expressionis ${}^{\beta}\!\surd\theta'=p+iq$, ita ut $P,Q,p,q$ sint reales, constat, si quantitates positivae $E,e$ angulique $F,f$ ita determinentur, ut sit $P=E\cos F$, $Q=E\sin F$, $p=e\cos f$, $q=e\sin f$, fore $e={}^{\beta}\!\surd E$, atque $f$ aequalem alicui ex angulis
\[
\frac1\beta F,\quad\frac1\beta(F+360^0),\quad\frac1\beta(F+720^0)\dotsfour\frac1\beta(F+(\beta-1)360^0)
\]
Determinabitur itaque $f$\textsuperscript{.}\ ger sectionem anguli $F$ in $\beta$ partes, at extractione radicis ${}^{\beta}\!\surd E$ sequenti modo supersedere possumus. Quodvis productum $r^kR^K$ partem suam realem habet communem cum $r^{-k}R^{-K}$, partes imaginariae autem factorem $i$ implicantes in his productis aequales sed oppositae erunt. Hinc sponte sequitur $[r^{-1},R^{-1}]=p-iq=e(\cos f-i\sin f)$, adeoque
\[
[r,R]\times[r^{-1},R^{-1}]=e^2
\]

Sed productum illud per theorema art.\ 4 fit
\[
\begin{aligned}
&{}=[1,1]+R[r^{g^\alpha-1},1]+R^2[r^{g^{2\alpha}-1},1]+\text{etc.}\\
&\qquad{}+R^{\beta\gamma-1}[r^{g^{\alpha\beta\gamma-\alpha}-1},1]\\
&{}=\beta\gamma+R(\beta\gamma,g^\alpha-1)+R^2(\beta\gamma,g^{2\alpha}-1)+\text{etc.}\\
&\qquad{}+R^{\beta\gamma-1}(\beta\gamma,g^{\alpha\beta\gamma-\alpha}-1)
\end{aligned}
\]
quae quantitas determinabilis est, si $R$ omnesque periodi $\beta\gamma$ terminorum cognitae supponuntur. Determinatio ipsius $e$ itaque solam extractionem radicis quadratae postulat.

In casu speciali, ubi $\alpha=1$, singulae periodi $(\beta\gamma,g^\alpha-1)$, $(\beta\gamma,g^{2\alpha}-1)$ etc.\ manifesto sunt $=r+r^2+r^3+r^4+\text{etc.}+r^{n-1}$, adeoque
\[
\begin{aligned}
ee&=\beta\gamma+(R+R^2+R^3+\text{etc.}+R^{\beta\gamma-1})(r+r^2+r^3+\text{etc.}+r^{n-1})\\
&=\beta\gamma+1=n
\end{aligned}
\]
\SourceLeaf{261}
siquidem $r$ et $R$ radices ab $1$ diversas exhibere supponuntur, et proin semper $e=\surd n$ (Disq. arithm. art. 360 fin.).

\Article{10}
Hactenus disquisitionem nostram summa generalitate instituimus, ut valores quoscunque numerorum $\alpha,\beta,\gamma$ complectatur: abhinc vero ad casum magis limitatum, ubi $\alpha=1$, transibimus, qui ad disquisitiones foecundissimas et elegantissimas viam nobis sternet. Exprimet itaque signum $[r,R]$ functionem
\[
r+Rr^g+R^2r^{g^2}+R^3r^{g^3}+\text{etc.}+R^{n-2}r^{g^{n-2}}
\]
ubi $n$ est numerus primus, $r$ indefinite radix aequationis $x^n-1=0$ (radice $1$ non excepta), $R$ indefinite radix aequationis $x^\beta-1=0$, denotante $\beta$ divisorem datum ipsius $n-1$, denique $g$ integer, qui est radix primitiva determinata pro modulo $n$. Porro brevitatis caussa scribemus
\[
\begin{aligned}
1+r+r^2+r^3+\text{etc.}+r^{n-1}&=s\\
1+R+R^2+R^3+\text{etc.}+R^{n-2}&=S
\end{aligned}
\]
unde patet $s$ fieri $=n$ pro $r=1$, sed $s=0$ pro quovis alio valore ipsius $r$, et perinde $S=n-1$ pro $R=1$, sed $S=0$ pro quovis alio valore ipsius $R$.

Per art. 3 itaque habemus $[1,R]=S$, $[r,1]=s-1$; porro pro quovis valore integri $m$ per $n$ non divisibili $[r^m,R]=R^{-\operatorname{ind}m}[r,R]$, aut generalius $[r^m,R^M]=R^{-M\operatorname{ind}m}[r,R^M]$, ubi $\operatorname{ind}m$ est exponens potestatis numeri $g$ secundum modulum $n$ ipsi $m$ congruae. Applicando hanc transformationem ad ea, quae in art. 5 docuimus, sequitur, productum e duabus pluribusve functionibus talibus $[r^h,R^H]$ reduci ad formam hanc
\[
A[1,R^\lambda]+B[r,R^\lambda]
\]
ubi $A$ et $B$ erunt functiones rationales ipsius $R$ cum coëfficientibus integris, atque $\lambda$ aggregatum omnium valorum ipsius $H$. Magni momenti erit, huiusmodi transformationes ad algorithmum expeditum reducere, ad quem finem imprimis indoles producti e duabus functionibus propius nobis consideranda erit.

\Article{11}
Productum $[r,R^\mu]\times[r,R^\nu]$ per theorema art. 4 fit $=$
\SourceLeaf{262}
\[
\begin{aligned}
&[r^2,R^{\mu+\nu}]+R^\mu[r^{g+1},R^{\mu+\nu}]+R^{2\mu}[r^{g^2+1},R^{\mu+\nu}]\\
&\qquad+R^{3\mu}[r^{g^3+1},R^{\mu+\nu}]+\text{etc.}+R^{(n-2)\mu}[r^{g^{n-2}+1},R^{\mu+\nu}]
\end{aligned}
\]
Inter $n-1$ exponentes $2,g+1,g^2+1,g^3+1$ etc. $g^{n-2}+1$ unus tantum reperitur per $n$ divisibilis, puta $g^{\frac12(n-1)}+1$, aggregati itaque nostri terminus respondens erit $R^{\frac12(n-1)\mu}[1,R^{\mu+\nu}]$: hic terminus erit $=0$, quoties non est $R^{\mu+\nu}=1$, et $=(n-1)R^{\frac12(n-1)\mu}=\pm(n-1)$, pro $R^{\mu+\nu}=1$. Partes reliquae aggregati nostri, quarum summam statuemus $=\Omega$, sequenti modo transformantur:
\[
\begin{aligned}
[r^2,R^{\mu+\nu}]&=R^{-(\mu+\nu)\operatorname{ind}2}[r,R^{\mu+\nu}]\\
R^\mu[r^{g+1},R^{\mu+\nu}]&=R^{\mu-(\mu+\nu)\operatorname{ind}(g+1)}[r,R^{\mu+\nu}]\\
R^{2\mu}[r^{g^2+1},R^{\mu+\nu}]&=R^{2\mu-(\mu+\nu)\operatorname{ind}(g^2+1)}[r,R^{\mu+\nu}]\\
R^{3\mu}[r^{g^3+1},R^{\mu+\nu}]&=R^{3\mu-(\mu+\nu)\operatorname{ind}(g^3+1)}[r.R^{\mu+\nu}]\\
&\text{etc.}
\end{aligned}
\]

Hinc colligimus
\[
\text{I.}\qquad\Omega=[r,R^{\mu+\nu}]\times\Sigma R^{\mu\operatorname{ind}x-(\mu+\nu)\operatorname{ind}(x+1)}
\]
si pro $x$ successive substituantur valores $1,g,g^2,g^3\dotsfour g^{n-2}$ excepto hoc $g^{\frac12(n-1)}$, seu quod manifesto eodem redit, si pro $x$ substituantur valores $1,2,3,4\dotsfour n-2$, quoniam valores hi illis (etsi ordine mutato) congrui sunt secundum modulum $n$.

Statuendo integro $y$ ipsi $x$ reciproco secundum modulum $n$, i. e. ita determinatum, ut fiat $xy\equiv1\pmod n$, erit $\operatorname{ind}x\equiv-\operatorname{ind}y\pmod{n-1}$, atque $\operatorname{ind}(x+1)+\operatorname{ind}y\equiv\operatorname{ind}(xy+y)\equiv\operatorname{ind}(1+y)\pmod{n-1}$; hinc fit
\[
\begin{aligned}
\mu\operatorname{ind}x-(\mu+\nu)\operatorname{ind}(x+1)&\equiv-\mu\operatorname{ind}y-(\mu+\nu)\{\operatorname{ind}(y+1)-\operatorname{ind}y\}\\
&\equiv\nu\operatorname{ind}y-(\mu+\nu)\operatorname{ind}(y+1)
\end{aligned}
\]
Quamobrem quum numeri ipsis $1,2,3\dots n-2$ reciproci cum his ipsis ordine tantum mutato conveniant, etiam erit
\[
\text{II.}\qquad\Omega=[r,R^{\mu+\nu}]\times\Sigma R^{\nu\operatorname{ind}y-(\mu+\nu)\operatorname{ind}(y+1)}
\]
substituendo pro $y$ successive numeros $1,2,3\dots n-2$. Eadem formula immediate ex I derivatur, quum manifesto numeros $\mu,\nu$ inter se permutare liceat.

Denique statuendo integrum $z$ ipsi $x+1$ reciprocum secundum 
\SourceLeaf{263}
modulum $n$, sive $xz+z\equiv1\pmod n$, erit $\operatorname{ind}(1-z)\equiv\operatorname{ind}x+\operatorname{ind}z\pmod{n-1}$, $\operatorname{ind}(x+1)\equiv-\operatorname{ind}z\pmod{n-1}$ adeoque
\[
\begin{aligned}
\mu\operatorname{ind}x-(\mu+\nu)\operatorname{ind}(x+1)&\equiv\mu(\operatorname{ind}(1-z)-\operatorname{ind}z)+(\mu+\nu)\operatorname{ind}z\\
&\equiv\mu\operatorname{ind}(1-z)+\nu\operatorname{ind}z
\end{aligned}
\]
Quare quum percurrente $x$ valores $1,2,3\dots n-2$, numerus $z$ percurrere debeat valores $2,3,4\dots n-1$ (etsi alio ordine), nanciscimur expressionem tertiam
\[
\text{III.}\qquad\Omega=[r,R^{\mu+\nu}]\times\Sigma R^{\mu\operatorname{ind}(1-z)+\nu\operatorname{ind}z}
\]
substituendo pro $z$ successive valores $2,3,4\dots n-1$, aut si mavis
\[
\begin{aligned}
\text{IV.}\qquad\Omega&=[r,R^{\mu+\nu}]\times\Sigma R^{\mu\operatorname{ind}(n+1-z)+\nu\operatorname{ind}z}\\
&=[r,R^{\mu+\nu}]\times\Sigma R^{\mu\operatorname{ind}z+\nu\operatorname{ind}(n+1-z)}
\end{aligned}
\]
Quum habeatur $\operatorname{ind}(1-z)=\frac12(n-1)+\operatorname{ind}(z-1)$, productum nostrum ita quoque exhiberi poterit:
\[
\begin{aligned}
[r,R^\mu]\times[r,R^\nu]&=R^{\frac12(n-1)\mu}\{[1,R^{\mu+\nu}]+[r,R^{\mu+\nu}]\times\Sigma R^{\mu\operatorname{ind}(z-1)+\nu\operatorname{ind}z}\}\\
&=R^{\frac12(n-1)\nu}\{[1,R^{\mu+\nu}]+[r,R^{\mu+\nu}]\times\Sigma R^{\mu\operatorname{ind}z+\nu\operatorname{ind}(z-1)}\}
\end{aligned}
\]
ubi semper pro $z$ substituendi concipiuntur valores $2,3,4\dotsfour n-1$.

Ceterum in omnibus his formulis pro numeris
\[
\mu\operatorname{ind}x-(\mu+\nu)\operatorname{ind}(x+1),\quad\nu\operatorname{ind}y-(\mu+\nu)\operatorname{ind}(y+1),\quad\mu\operatorname{ind}(1-z)+\nu\operatorname{ind}z
\]
etc. manifesto ipsorum residua minima secundum modulum $\beta$ substitui poterunt.

Si $\mu+\nu\equiv0\pmod\beta$ erit
\[
\begin{aligned}
[r,R^\mu][r,R^\nu]&=(n-1)R^{\frac12(n-1)\mu}\\
&\quad+(r+r^2+r^3+\dots+r^{n-1})\\
&\qquad\times(1+R^\mu+R^{2\mu}+R^{3\mu}+\dots+R^{(n-2)\mu}-R^{\frac12(n-1)\mu})
\end{aligned}
\]

\Article{12}
Productum $[1,R^\mu]\times[r,R^\nu]$ per theorema art. 4 fit
\[
\begin{aligned}
&=[r,R^{\mu+\nu}]+R^\mu[r,R^{\mu+\nu}]+R^{2\mu}[r,R^{\mu+\nu}]+\text{etc.}+R^{(n-2)\mu}[r,R^{\mu+\nu}]\\
&=[r,R^{\mu+\nu}]\times(1+R^\mu+R^{2\mu}+R^{3\mu}+\text{etc.}+R^{(n-2)\mu})\\
&=[r,R^{\mu+\nu}]\times\frac{n-1}\beta(1+R^\mu+R^{2\mu}+R^{3\mu}+\text{etc.}+R^{(\beta-1)\mu})
\end{aligned}
\]

\SourceLeaf{264}
Hinc productum $[1,R^\mu]\times[1,R^\nu]$ evolvitur in
\[
\frac{n-1}\beta[1,R^{\mu+\nu}]\times(1+R^\mu+R^{2\mu}+R^{3\mu}+\text{etc.}+R^{(\beta-1)\mu})
\]
Nullo iam negotio generaliter productum $[r^m,R^\mu]\pdot[r^{m'},R^{\mu'}]$ erui poterit, quum enim fiat $[r^m,R^\mu]=R^{-\mu\operatorname{ind}m}[r,R^\mu]$ pro valore ipsius $m$ per $n$ non divisibili, et $=[1,R^\mu]$ pro valore divisibili, et quum similis transformatio de factore altero $[n^{m'},R^{\mu'}]$ valeat, multiplicatio vel ad problema art. praec. reducetur, vel ad casus eos, quos in hoc art. consideravimus.
